\chapter{Implementação e Avaliação Experimental}

Este capítulo apresenta a implementação do sistema desenvolvido e os experimentos conduzidos com o objetivo de avaliar seu desempenho, escalabilidade e corretude geométrica. Em particular, busca-se analisar o impacto de diferentes estratégias nas etapas de filtragem e refinamento, bem como comparar os algoritmos SAT e GJK/EPA em um ambiente de execução web, no qual restrições de desempenho e responsividade são relevantes.

Os experimentos foram realizados em um computador com as seguintes especificações:
\begin{itemize}
	\item Processador: AMD Ryzen 5 5500U
	\item Placa de Vídeo: AMD Radeon Vega 7 (integrada)
	\item Memória RAM: 24 GB DDR4
	\item Sistema Operacional: Ubuntu 24.04 LTS
\end{itemize}

\section{Arquitetura e Implementação}

Foi desenvolvida uma biblioteca em \textit{TypeScript} responsável pela implementação dos algoritmos de detecção de colisão, sendo reutilizada tanto na aplicação de visualização quanto no ambiente de testes. A engine permite a execução em dois modos distintos: utilizando o algoritmo do Teorema do Eixo Separador (SAT) ou a combinação dos algoritmos GJK/EPA.

A arquitetura do sistema foi projetada de forma modular, separando a lógica geométrica, os algoritmos de detecção e resposta à colisão e a camada de visualização, delegando esta última à aplicação cliente. Essa separação favorece a reutilização de código e possibilita a comparação direta entre diferentes abordagens de detecção de colisão.

Os objetos convexos são representados como listas de vértices ordenados, associadas a um conjunto de restrições. Cada restrição atua sobre um par de partículas, aplicando um processo de relaxamento que permite modelar tanto estruturas rígidas quanto deformáveis. Essa modelagem possibilita a integração da detecção de colisão com simulações físicas baseadas em partículas. 

Este trabalho disponibiliza uma implementação pública com o objetivo de permitir a reprodução dos experimentos e a extensão das técnicas apresentadas~\footnote{Disponível em: \url{https://github.com/devdiegomatos/simplified-physics-animation-web}}.

\section{Ambiente de Visualização}

Com o objetivo de permitir a análise qualitativa do comportamento dos algoritmos, foi desenvolvido um site em \textit{Vue.js 3} com duas \textit{demos}~\footnote{https://devdiegomatos.github.io/simplified-physics-animation-web/}, utilizando \textit{p5.js} para renderização do canvas. Vale ressaltar que a biblioteca é independente da visualização, as escolhas de \textit{Vue.js 3} e \textit{p5.js} foram apenas preferência do autor.

A primeira demonstração consiste em um ambiente contendo múltiplos objetos convexos, com número de vértices variando de 3 e 8, inseridos em uma caixa com obstáculos. Esse cenário permite a observação da dinâmica de colisões em tempo real. A interface possibilita o controle de parâmetros relevantes, como a quantidade de objetos, o algoritmo de detecção de colisão (SAT ou GJK/EPA) e a estratégia de filtragem, incluindo o uso de uma grade uniforme. A Figura~\ref{fig:demo1} ilustra esse ambiente.

\begin{figure}[htb]
    \centering
    % Primeira minipage (Imagem A)
    \begin{minipage}{0.48\textwidth}
        \centering
        \includegraphics[width=\textwidth]{demo1.png}
        \caption{Fonte: Autor.}
        \label{fig:demo1}
    \end{minipage}
    \hfill % Espaço entre as imagens
    % Segunda minipage (Imagem B)
    \begin{minipage}{0.48\textwidth}
        \centering
        \includegraphics[width=\textwidth]{demo2.png}
        \caption{Fonte: Autor.}
        \label{fig:demo2}
    \end{minipage}
    
    % \caption{Legenda principal para ambas as imagens}
    % \label{fig:conjunto}
\end{figure}

A segunda demonstração apresenta simulações de tecidos e treliças, evidenciando a interação entre objetos deformáveis. Nesse cenário, as restrições entre partículas são responsáveis pela manutenção estrutural, enquanto o sistema de colisão garante a interação física entre os elementos. A Figura~\ref{fig:demo2} apresenta um exemplo desse ambiente.

\section{Metodologia de Testes}

Os experimentos foram conduzidos com quatro objetivos principais:
\begin{itemize}
    \item Avaliar o impacto da utilização de uma estrutura de filtragem no desempenho do sistema;
    \item Comparar o desempenho dos algoritmos SAT e GJK/EPA na fase de refinamento;
    \item Analisar a escalabilidade das abordagens em função do número de objetos.
    \item Comparar o desempenho da implementação desenvolvida com uma biblioteca comumente usada em jogos. Aqui, usamos a  biblioteca \textit{Planck.js}~\cite{planckjs}.
\end{itemize}
Os testes foram realizados sem a etapa de visualização, a fim de isolar o custo computacional da detecção de colisão. Para garantir a reprodutibilidade dos experimentos, foi utilizada geração pseudoaleatória com semente fixa por meio da biblioteca \textit{seedrandom}~\cite{bau2019seedrandom}, uma vez que a função \textit{Math.random} não permite controle de semente.

O cenário de teste corresponde ao ambiente apresentado na Figura~\ref{fig:demo1}, no qual objetos convexos são posicionados aleatoriamente dentro de uma região delimitada. Foram variados o número de objetos e o tamanho do espaço. A geração das posições iniciais segue uma distribuição de Poisson (\textit{Poisson Disk Sampling}), que tem como característica gerar amostras espacialmente uniformes com distância mínima garantida entre pontos e melhorando a estabilidade inicial da simulação~\cite{bridson2007poisson}. Essa propriedade evita a concentração de objetos em regiões específicas e reduz a ocorrência de interpenetrações iniciais, contribuindo para a construção de cenários mais homogêneos e adequados à avaliação do desempenho dos algoritmos de detecção de colisão.

Foram considerados quatro casos experimentais:
\begin{itemize}
    \item SAT sem fase de filtragem;
    \item SAT com grade uniforme na fase de filtragem;
    \item GJK/EPA sem fase de filtragem;
    \item GJK/EPA com grade uniforme na fase de filtragem.
\end{itemize}
O número de objetos foi variado de 100 a 1000, em incrementos de 100. Para cada configuração, foram realizadas 180 quadros independentes, sendo computadas as médias das métricas coletadas.

As métricas analisadas incluem o tempo médio de execução, o número de testes de colisão realizados e tempo decorrido em cada etapa da simulação, permitindo avaliar tanto o custo computacional quanto a corretude dos algoritmos. Essa abordagem permite a visualização simultânea de todas as configurações, preservando a comparabilidade entre os métodos.

A biblioteca \textit{Planck.js} é um motor de física 2D para JavaScript, reescrito a partir do Box2D, amplamente utilizado no desenvolvimento de jogos digitais. A \textit{Planck.js} emprega um método baseado em impulso e usa uma estrutura de árvore hierárquica dinâmica de AABB na fase de filtragem, enquanto a detecção de colisão na fase de refinamento baseia-se no GJK e EPA.

Para garantir a comparabilidade dos resultados, foram utilizados cenários equivalentes aos descritos anteriormente, mantendo-se consistentes os parâmetros de número de objetos, distribuição espacial e número de execuções. Devido à forma como a biblioteca \textit{Planck.js} é implementada, foi analisada apenas a métrica de tempo médio de execução.

\section{Resultados Experimentais}

Os resultados obtidos permitem analisar o comportamento dos algoritmos sob diferentes configurações e validar empiricamente as hipóteses estabelecidas.

\begin{figure}[htb]
    \centering
    \includegraphics[width=0.85\textwidth]{comparativo-testes-colisao.png}
    \caption{Número médio de testes de colisão em função do número de objetos, com e sem utilização de grade uniforme na filtragem, da nossa implementação. Fonte: Autor.}
    \label{fig:testes-comparativo}
\end{figure}

A Figura~\ref{fig:testes-comparativo} evidencia o impacto direto da etapa da filtragem na redução do número de testes de colisão. A utilização da grade uniforme reduz drasticamente a quantidade de pares avaliados, refletindo diretamente na melhoria de desempenho observada. O algoritmo de refinamento não influencia este teste, pois a etapa de filtragem é independente. Devido à diferença de ordens de grandeza entre os cenários com e sem fase de filtragem, os resultados são apresentados em escala logarítmica no eixo vertical.

\begin{figure}[htb]
    \centering
    \includegraphics[width=0.85\textwidth]{profiling-sat.png}
    \caption{Tempo relativo por etapas da nossa implementação no modo SAT. Fonte: Autor.}
    \label{fig:profiling-sat}
\end{figure}

\begin{figure}[htb]
    \centering
    \includegraphics[width=0.85\textwidth]{profiling-gjk.png}
    \caption{Tempo relativo por etapas da nossa implementação no modo GJK. Fonte: Autor.}
    \label{fig:profiling-gjk}
\end{figure}

% A análise da Figura~\ref{fig:profiling} revela um alto impacto da etapa de relaxamento da nossa implementação, o que pode ser atribuído à forma como as restrições são processadas e à necessidade de iterar sobre os pares de partículas para aplicar as correções. Essa etapa é fundamental para garantir a estabilidade da simulação, mas também representa um custo computacional significativo, especialmente em cenários com um grande número de objetos e restrições. 

As Figuras~\ref{fig:profiling-sat} e \ref{fig:profiling-gjk} apresentam um gráfico de área empilhado onde o primeiro item da legenda está na base da pilha. Neste gráfico foi mensurado o tempo relativo de cada etapa da implementação usando SAT e GJK. A análise evidencia que a etapa de relaxamento domina o custo total de execução em todos os cenários. Esse comportamento decorre do método proposto por Jakobsen, que, embora simples e estável, requer múltiplas iterações para convergir, aumentando o custo proporcionalmente ao número de restrições e de passos de relaxamento. Neste trabalho foram usados 5 iterações de relaxamento. Observa-se também que a fase de filtragem e refinamento são, respectivamente, as mais custosas depois do relaxamento, com o GJK apresentando maior custo.

Esses resultados indicam que ganhos de desempenho podem ser obtidos tanto pela redução do número de iterações de relaxamento e o uso de outras estruturas de dados na fase de filtragem, como estruturas hierárquicas. Adicionalmente, a natureza paralelizável dessas etapas sugere o uso de execução multi-thread ou GPU como estratégias promissoras. 

\begin{figure}[htb]
    \centering
    \includegraphics[width=0.85\textwidth]{comparativo-tempo-execução.png}
    \caption{Tempo médio de execução em função do número de objetos para os diferentes algoritmos de refinamento comparado à biblioteca \textit{Planck.js}. Fonte: Autor.}
    \label{fig:tempo-comparativo}
\end{figure}

Na Figura~\ref{fig:tempo-comparativo}, observa-se que a biblioteca \textit{Planck.js} apresenta desempenho superior em todos os cenários avaliados. Esse comportamento pode ser atribuído a diferenças fundamentais na abordagem de simulação e detecção de colisão. Enquanto a implementação proposta requer múltiplas iterações de relaxamento para convergência, a \textit{Planck.js} adota um modelo baseado em impulsos, no qual a resolução de colisões é realizada diretamente no nível dos corpos rígidos, acumulando e aplicando forças resultantes de forma mais eficiente. Outro fator relevante é o uso de estruturas hierárquicas na fase de filtragem, como árvores dinâmicas de volumes limitadores.
